\documentclass[12pt]{article}
\usepackage[utf8]{inputenc}

\usepackage{mystyle}

\title{Stability Conditions}
\author{Joseph Sullivan}
\date{October 2024}

\tikzset{
    wb/.style={
        ->,preaction={draw, line width=5pt, white}
    }
}

\DeclareMathOperator{\Kom}{Kom}
\DeclareMathOperator{\Cyl}{Cyl}
\DeclareMathOperator{\Stab}{Stab}
\DeclareMathOperator{\Slice}{Slice}
\DeclareMathOperator{\chern}{ch}
\DeclareMathOperator{\todd}{td}

\newcommand{\cQ}{\mathcal{Q}}
\newcommand{\cZ}{\mathcal{Z}}

\newcommand{\rddots}{\text{\reflectbox{$\ddots$}}}


\begin{document}

\maketitle
\tableofcontents

\section{Algebraic t-structures}

\begin{defn}
    Let $\cD$ be a $k$-linear triangulated category, and $X,Y\in \cD$. Define
    \[\Hom^\bullet(X,Y) = \bigoplus_{i\in \Z} \Hom(X,Y[i]),\]
    and $\End^\bullet(X)$ similarly (which is a graded ring using shifts as necessary to compose).

    \begin{itemize}
        \item We say $E$ is \textbf{exceptional} if $\End^\bullet(E) = k\cdot 1_E$.
        
        \item We say $E_0,\dots,E_n$ (order matters!) is an \textbf{exceptional collection} if each $E_i$ is exceptional and $\Hom^\bullet(E_j,E_i)=0$ when $j>i$ (no backwards Homs).
    \end{itemize}

    Now, assume $E_0,\dots,E_n$ is an exceptional collection, and we say the collection is
    \begin{itemize}
        \item \textbf{strong} if $\Hom^\bullet(E_i,E_j) = \Hom^0(E_i,E_j)$ (remaining Homs are concentrated in degree 0).

        \item \textbf{Ext} if $\Hom^\bullet(E_i,E_j) = \Hom^{>0}(E_i,E_j)$ when $i\neq j$ (remaining Homs are in positive degree).

        \item \textbf{full} if the $E_i$ generate $\cD$ as a triangulated category (the smallest triangulated category containing is the $E_i$ is $\cD$, so generated under cones/shifts).
    \end{itemize}
\end{defn}

The following is due to \cite{macri2007stabilityconditionscurves}.
\begin{prop} \label{algebraic-t-structure}
    Let $E_0,\dots,E_n$ be a full Ext exceptional collection. Then, $\<E_0,\dots,E_n\>$ is the heart of a bounded t-structure. We refer to this as an \textbf{algebraic t-structure}.
\end{prop}

\begin{rmk}
    These abelian categories are \textit{finite length}, i.e., every object has a finite Jordan-Holder decomposition series, because any
    \[X \in \<E_0,\dots,E_n\>\]
    is obtained from a finite sequence of extensions of the $E_i$. {\color{red} The simple objects are exactly the $E_i$. (this is true, right??)}
\end{rmk}

\begin{eg}
    Beilinson's resolution of the diagonal tells us that
    \[\cO(-n),\dots,\cO\]
    and
    \[\cO,\Omega^1(1),\dots,\Omega^n(n)\]
    are both full strong exceptional collections. In particular,
    \begin{align*}
        \cO(-n)[n],\dots,\cO[0]\\
        \cO[0],\dots,\Omega^n(n)[-n]
    \end{align*}
    are then Ext exceptional collections. Now, this lemma tells us that
    \[\<\cO(-n)[n],\dots,\cO[0]\>\]
    is a heart of a bounded t-structure. In fact, 
    \[\<\cO(-n)[n],\dots,\cO[0]\> = \{0\to V^{-n} \otimes_k \cO(-n) \to \cdots \to V^0 \otimes_k \cO \to 0\},\]
    because this contains each of these objects and is extension closed (an extension of $A$ by $B$ is a cone over a morphism $A[-1] \to B$, and Beilinson gets us that these are honest chain maps, and such cones are still in $\cA_W$).
\end{eg}

\begin{defn}
    Denote the full additive subcategory of $D^b(\P^n)$
    \[\cA_W = \<\cO(-n)[n],\dots,\cO[0]\>\]
    the \textbf{category of linear complexes}.
\end{defn}

\begin{rmk}
    {\color{red} I believe that there should be some sort of Koszul duality reason that if you take the full strong exceptional collection $\cO,\Omega^1(1),\dots,\Omega^n(n)$, take their direct sum to get a tilting sheaf, and use Bondal's result to get a derived equivalence to $\rmod A$ (for $A$ the endomorphism ring of the tilting sheaf), the equivalence will send $\rmod A$ exactly to $\cA_W$. It seems to me that when we use the opposite exceptional collection we get the opposite linear complexes}.
\end{rmk}

\section{Torsion Pair t-structures}

One large source of examples of t-structures is by what is called ``tilting at a torsion pair''.

\begin{defn}
    Let $\cA$ be an abelian category. A \textbf{torsion pair} $(\cT,\cF)$ is a pair of (strictly) full subcategories of $\cA$ satisfying
    \begin{enumerate}[(i)]
        \item $\Hom(\cT,\cF) = 0$
        \item Every object $A \in \cA$ fits in a short exact sequence
        \[0\to T \to A \to F \to 0\]
        for some $T\in \cT$ and $F\in \cF$.
    \end{enumerate}
\end{defn}

Now, we explain how to tilt at a torsion pair.

\begin{thmdef}
    Let $\cD$ be a triangulated category, and let $\cA$ be the heart of a bounded t-structure on $\cD$. Let $(\cT,\cF)$ be a torsion pair on $\cA$. Then, the following is the heart of a bounded t-structure on $\cD$:
    \[\cA^\sharp = \<\cF[1],\cT\> = \{X\in \cD \mid H_\cA^0(X) \in \cT, H_\cA^{-1}(X) \in \cF, H_\cA^i(X)=0 \text{ otherwise}\}.\]
\end{thmdef}

\begin{proof}
    We give a sketch.

    First, $\Hom(\cA^\sharp,\cA^\sharp[i])=0$ for $i<0$ because
    \[\Hom(\cF[1],\cF[1+i])=\Hom(\cF[1],\cT[i])=\Hom(\cT,\cT[i])=\Hom(\cT,\cF[1+i])=0,\]
    (where the last equality comes from the definition of a torsion pair) imply
    \[\Hom(\<\cF[1],\cT\>, \<\cF[1],\cT\>[i]) = 0.\]

    Now, we need to produce t-structure cohomology filtrations for $\cA^\sharp$. Let $E \in \cD$. Our bounded t-structure $\cA$ gives us a filtration
    \begin{cd}[column sep=0.25cm]
        0=E^{-n} \ar[rr] && E^{-n+1} \ar[dl] \rar & \cdots \rar & E^{-1} \ar[rr] && E^0 \rar \ar[dl] & \cdots \rar & E^{n-1} \ar[rr] && E^n=E \ar[dl]\\
        & A^{-n} \ar[ul,dashed] && && A^0 \ar[ul,dashed] && && A^n \ar[ul,dashed] &
    \end{cd}
    where $A^i \in \cA[-i]$. Then, we ``refine'' the filtration using the torsion pair as follows. By the torsion pair, we have a short exact sequence
    \[0 \to T^i \to A^i \to F^i \to 0,\]
    which will give us a distinguished triangle in $\cA$. Thinking of $E^i$ as an extension of $A^i$ by $E^{i-1}$, we will lift come up with $\tilde{T}^i$ which is an extension of $T^i$ by $E^{i-1}$, such that we have distinguished triangles.
    \begin{cd}[column sep=0.25cm]
        E^{i-1} \ar[rr] && \tilde{T}^i \ar[ld]\ar[rr] && E^i \ar[ld]\\
        & T^i \ar[lu,dashed] && F^i \ar[lu,dashed] &
    \end{cd}
    We can do this by defining
    \[\tilde{T}^i = C(T^i[-1] \to E^{i-1}),\]
    and then getting a map $\tilde{T}^i \to E^i$ by pulling back $E^{i-1}\to E^i$, after noticing that the composition $T^i[-1] \to E^{i-1} \to \tilde{T}^i$ is 0 (and using that $\Hom$ is a cohomological functor). Then, the cone over $E^{i-1}\to \tilde{T}^i$ is $T^i$ by definition, and the cone over $\tilde{T}^i \to E^i$ is $F^i$ by the octahedral axiom (i.e., analyze the composition $E^{i-1}\to \tilde{T}^i \to E^i$).
    
    The point now is to filter $E$ by the $\tilde{T}^i$ instead of the $E^i$.
    
    Putting two of these diagrams next to each other and deleting the left/right ends, we get the following.
    \begin{cd}[column sep=0.25cm]
        \tilde{T}^{i-1} \ar[rr] && E^{i-1} \ar[ld] \ar[rr] && \tilde{T}^i \ar[ld]\\
        & F^{i-1} \ar[lu,dashed] && T^i \ar[lu,dashed] &
    \end{cd}
    We now need to compute the cone over $\tilde{T}^{i-1} \to \tilde{T}^i$, which we denote by $\tilde{A}^i$. Applying the octahedral axiom, we get a distinguished triangle
    \[F^{i-1} \to \tilde{A}^i \to T^i \to F^{i-1}[1].\]
    Taking $\cA$ cohomology, and recognizing that
    \[H_\cA^j(F^{i-1}) = \begin{cases}
        F^{i-1}[i-1] & j=i-1\\
        0 & \text{otherwise}
    \end{cases}, \qquad H_\cA^j(T^i) = \begin{cases}
        T^i[i] & j=i\\
        0 & \text{otherwise}
    \end{cases}\]
    we then get exact sequences
    \[0 \to F^{i-1}[i-1] \to H_\cA^{i-1}(\tilde{A}^i) \to 0\]
    and
    \[0 \to H_\cA^i(\tilde{A}^i) \to T^i[i] \to 0,\]
    so in other words
    \begin{align*}
        H_\cA^{-1}(\tilde{A}^i[i]) &= F^{i-1}[i-1] \in \cF\\
        H_\cA^0(\tilde{A}^i[i]) &= T^i[i] \in \cT,
    \end{align*}
    so indeed
    \[\tilde{A}^i \in \cA^\sharp[-i],\]
    so we get an $\cA^\sharp$ t-structure cohomology filtration as desired.
\end{proof}

Tilting provides tons of examples of t-structures---starting with the standard t-structure $\cA \subseteq D^b(\cA)$, we can tilt over and over again to find more and more t-structures. One question we can ask is whether two given t-structures are tilts of each other. We can the following proposition to help us answer this questions.

\begin{prop}
    Let $\cA^\sharp, \cA$ be hearts of bounded t-structures on a triangulated category $\cD$, and assume
    \[\cA^\sharp \subseteq \<\cA[1],\cA\>.\]
    Then, $\cA^\sharp$ is a tilt of $\cA$ at the torsion pair
    \[\cT = \cA \cap \cA^\sharp, \qquad \cF = \cA \cap \cA^\sharp[-1].\]
\end{prop}

\begin{proof}
    It's clear that 
    \[\Hom(\cT,\cF) = 0,\]
    since $\cT \subseteq \cA^\sharp$ and $\cF\subseteq \cA^\sharp[-1]$.

    Now, we need to show that every object $A\in \cA$ is an extension of $\cF$ by $\cT$. To do this, we study the $\cA^\sharp$ cohomology filtration of $A$, which will be something like
    \begin{cd}[column sep=0.25cm]
        \cdots \ar[rr]&&A^{-2} \ar[rr] && A^{-1} \ar[dl]\ar[rr] && A^0 \ar[dl]\ar[rr] && A^1 \ar[dl]\ar[rr] && A^2 \ar[dl]\ar[rr] && \cdots\\
        && & H^{-1} \ar[ul,dashed] && H^0 \ar[ul,dashed] && H^1 \ar[ul,dashed] && H^2 \ar[ul,dashed] & &&
    \end{cd}
    and we want to argue that only $H^0,H^1$ are nonzero, so that we get a distinguished triangle
    \[H^0 \to A \to H^1 \to H^0[1],\]
    {\color{red} FINISH THIS!!}
\end{proof}

\begin{eg}
    Let $(X,H)$ be a polarized surface. We will define a tilt of $D^b(X)$ which will be used later to define a stability condition.

    Recall 
    
    Let $\mu \in (-\infty,\infty]$. Then, define
    \[\cT_\mu = \{\text{torsion sheaves}\} \cup \{\text{torsion-free sheaves with } \mu^- > \mu\}\]
    \[\cF_\mu = \{\text{torsion-free sheaves with $\mu^+ \leq \mu$}\}\]
\end{eg}



\section{Stability Condition}
Now that we have a heart, we can write down stability conditions. Recall that
\[\chern: K(\Coh(\P^n)) \to H^*(\P^n;\Q)\]
is an isomorphism onto its image. We claim that $\cO(-n),\dots,\cO$ is a basis for $K(\Coh(\P^n))$. To see this we can check that they have linearly independent chern characters,
\[K(\Coh(\P^n)) = \<\>\]


\nocite{*} % Include all references in refs.bib regardless of whether they are referenced or not
\bibliographystyle{plain} % We choose the "plain" reference style
\bibliography{2025spring/stability_on_proj_refs} % Entries are in the refs.bib file

\end{document}
